\hypertarget{references}{%
\chapter{References}\label{references}}

Works cited across Volume 1. Chapters cite by key ---
\texttt{{[}Anderson{]}}, \texttt{{[}Irvine{]}} --- and the keys resolve
here. Each work is entered once.

{[}Anderson{]} James P. Anderson, \emph{Computer Security Technology
Planning Study}, ESD-TR-73-51, Vol. I, Air Force Electronic Systems
Division, Hanscom AFB, Bedford, MA, October 1972. The reference monitor
and its three requirements are set out in §4.1.1.
\url{https://csrc.nist.gov/files/pubs/conference/1998/10/08/proceedings-of-the-21st-nissc-1998/final/docs/early-cs-papers/ande72.pdf}

{[}Irvine{]} Cynthia E. Irvine, ``The Reference Monitor Concept as a
Unifying Principle in Computer Security Education,'' in Proc. IFIP TC11
WG 11.8 First World Conference on Information Security Education, 1999,
pp.~27--37 (Naval Postgraduate School, June 1999).
\url{https://apps.dtic.mil/sti/tr/pdf/ADA423529.pdf}

{[}SaltzerSchroeder{]} Jerome H. Saltzer and Michael D. Schroeder, ``The
Protection of Information in Computer Systems,'' \emph{Proceedings of
the IEEE}, vol.~63, no. 9, pp.~1278--1308, September 1975. Origin of the
design principles, among them complete mediation, least privilege,
economy of mechanism, and fail-safe defaults.
\url{https://www.cs.virginia.edu/~evans/cs551/saltzer/}

{[}Schneider{]} Fred B. Schneider, ``Enforceable Security Policies,''
\emph{ACM Transactions on Information and System Security}, vol.~3, no.
1, pp.~30--50, February 2000. Characterizes the policies a monitor can
enforce by watching execution --- exactly the safety properties --- with
complete mediation and isolation among its conditions, and security
automata as the specification.
\url{https://www.cs.cornell.edu/fbs/publications/EnfSecPols.pdf}
